\newpage
\cleardoublepage
\chapter{Results}\label{chapter:results}

% TODO: lead-in paragraph stating the chapter's restructured weighting up
% front, mirroring the framing added to Chapter \ref{chapter:literature_review}.
% The headline results are (1) the practicability of synthetic training data,
% evaluated through the generator comparison of
% docs/synthetic-model-comparison (\Cref{sec:results_synthetic_data_practicability}),
% and (2) detection performance across the real/synthetic training-and-test
% regimes defined in docs/plans/2026-06-10_model-evaluation-strategy.md
% (\Cref{sec:results_regime_comparison}), since these are where the bulk of the
% thesis's actual experimental work went. Knowledge distillation and the
% embedded hardware check are demoted to a single smaller, explicitly-scoped-
% as-basic closing section (\Cref{sec:results_kd_and_embedded}), reflecting
% that only a single teacher/student comparison and a simple runtime check
% were carried out, not the fuller distillation-vs-direct-fine-tuning research
% program originally planned.

%%%%%%%%%%%%%%%%%%%%%%%%%%%%%%%%%%%%%%%%%%%%%%%%%
\section{Practicability of Synthetic Training Data: Generator Comparison}\label{sec:results_synthetic_data_practicability}
% TODO: reports the controlled comparison designed in
% docs/synthetic-model-comparison/01_experiment-design.md across the reduced
% class subset (docs/synthetic-model-comparison/02_class-selection.md) and the
% generator grid (docs/synthetic-model-comparison/03_api-models-landscape-and-pricing.md,
% 04_local-models-and-output-parameters.md): the incumbent
% \texttt{gemini-3.1-flash-image-preview}, other API models (gpt-image-2 at
% multiple tiers, Nano Banana 2 Lite/Pro), and local diffusion baselines
% (FLUX.1-schnell, RealVisXL+SDXL-Lightning, SD 3.5 Medium), each at full and
% prompt-length-compressed regimes where applicable
% (docs/synthetic-model-comparison/05_prompt-strategy-and-length-limits.md).
% Follows the three-axis methodology of
% docs/synthetic-model-comparison/06_evaluation-methodology.md.

\subsection{Qualitative Rubric Results}\label{sec:results_qualitative_rubric}
% TODO: blind, multi-rater scores per generator on the fixed rubric
% (anatomical correctness, diagnostic-feature fidelity, species identity,
% habitat plausibility, photorealism, artifact rate); report inter-rater
% agreement (Cohen's/Fleiss' kappa) alongside the scores, not as a footnote.

\subsection{Automatic Proxy Results}\label{sec:results_automatic_proxies}
% TODO: teacher-recognition confidence (SpeciesNet-on-synthetic top-1 and mean
% confidence) as the primary proxy, per-class FID/KID against real imagery,
% CLIP-score for prompt adherence, and MegaDetector auto-label yield per
% generator (a practicality signal, not just a quality one).

\subsection{Downstream Utility: Real-Test mAP per Generator}\label{sec:results_downstream_utility}
% TODO: the fixed detector/classifier trained on each cell's synthetic images,
% evaluated on the held-out real test set only (mixed and real-only
% breakouts per the evaluation strategy); the three-zebra-species fine-grained
% probe as the decisive look-alike test; per-generator €/image and
% images/hour alongside the accuracy numbers, following the single
% cross-axis ranking table sketched in
% docs/synthetic-model-comparison/06_evaluation-methodology.md \S7.

\subsection{Practicability Verdict}\label{sec:results_practicability_verdict}
% TODO: synthesize the three axes into a single practicability judgment --
% does the current generation of text-to-image models produce training data
% usable for this task at all, and was the production choice of Gemini
% justified against the alternatives on the axes that matter (real-test mAP
% and diagnostic fidelity), not merely on informal photorealism. This is the
% section that most directly answers the professor's original methodological
% objection (docs/synthetic-model-comparison/00_motivation-and-research-question.md).

%%%%%%%%%%%%%%%%%%%%%%%%%%%%%%%%%%%%%%%%%%%%%%%%%
\section{Detection Performance Across Training Regimes and Test Domains}\label{sec:results_regime_comparison}
% TODO: the thesis's other headline result, reporting the Tier 1/Tier 2
% structure of docs/plans/2026-06-10_model-evaluation-strategy.md \S9 for the
% YOLOv5s detector trained under the Band A-D data regimes
% (Methods \Cref{sec:band_structure}).

\subsection{Headline: Mixed Real+Synthetic mAP and the Real-Only Breakout}\label{sec:results_headline_mixed}
% TODO: Tier 1 table -- G=fine, D=mixed, B=all, full COCO-12 vector, alongside
% the real-only breakout (the public-comparison anchor) and the
% \texttt{mAP\_detect} class-agnostic analog (evaluation strategy \S7).

\subsection{Granularity Gap Decomposition}\label{sec:results_granularity_gaps}
% TODO: detect/coarse/fine mAP as a gap decomposition
% ($\Delta_{\text{coarse}}$, $\Delta_{\text{fine}}$, evaluation strategy \S4),
% backed by a TIDE error-type chart; interpret a large $\Delta_{\text{fine}}$
% with a small $\Delta_{\text{coarse}}$ as the expected look-alike-confusion
% signature for this dataset.

\subsection{Band $\times$ Granularity: Does Training Regime Change Accuracy?}\label{sec:results_band_granularity}
% TODO: the dataset's core scientific comparison -- the 4-band $\times$
% \{fine, coarse\} $\times$ \{mAP, mAP50\} grid (evaluation strategy \S6a),
% reported on the mixed domain with the real-only grid alongside it as the
% synthetic-vs-real-training tie-breaker, since Band A/B are partly scored
% on their own training domain under \texttt{mixed}.

\subsection{Domain-Shift Delta: Real vs.\ Synthetic Test Performance}\label{sec:results_domain_shift}
% TODO: the paired $\Delta_{\text{domain}}(\text{class}) = \text{mAP}_{\text{real}} -
% \text{mAP}_{\text{synthetic}}$ per class (evaluation strategy \S6b), reported
% as a distribution per band; this is also the watchdog check for the
% \texttt{mixed} default itself -- state explicitly whether the check passed
% (no revision of the default axes triggered) or whether a discrepancy emerged.

\subsection{Look-Alike Within-Group Confusion}\label{sec:results_lookalike_confusion}
% TODO: per-look-alike-group confusion table plus the block-diagonal confusion
% matrix restricted to the frozen look-alike groups
% (docs/plans/2026-06-11_lookalike-groups-review.md), using the same
% taxonomic-tolerance groups mined during ground-truth review
% (Methods \Cref{sec:gt_annotation_integrity}).

%%%%%%%%%%%%%%%%%%%%%%%%%%%%%%%%%%%%%%%%%%%%%%%%%
\section{Knowledge Distillation and Embedded Feasibility}\label{sec:results_kd_and_embedded}
% TODO: explicitly frame this section, in its opening paragraph, as
% intentionally reduced in scope relative to the original research question --
% a single basic comparison and a single runtime check, not an exploration of
% the wider design space surveyed in
% \Cref{sec:lit_knowledge_distillation,sec:lit_embedded_inference}. State the
% time-budget reason plainly rather than presenting the reduced scope as if it
% were the original plan.

\subsection{Distillation vs.\ Direct Fine-Tuning: A Single Basic Comparison}\label{sec:results_kd_basic_comparison}
% TODO: Setup A vs.\ Setup B (Methods \Cref{sec:augmentation_setups}), byte-identical
% augmentation pipelines, isolating the effect of the teacher signal alone;
% report the same headline mAP figure as \Cref{sec:results_headline_mixed} for
% both setups side by side. State clearly that this is one direct comparison
% at one teacher/student pairing, not a sweep over distillation strategies,
% temperature values, or teacher architectures.

\subsection{Runtime Benchmark on AX Visio Hardware}\label{sec:results_ax_visio_benchmark}
% TODO: the export pipeline used to get the trained YOLOv5s model onto the
% \textit{AX Visio} device, and the measured wall-clock inference
% latency/throughput (FPS) on-device; report it alongside the equivalent
% Raspberry Pi 5 proxy measurement to close the loop on the hardware-proxy
% reasoning of docs/2026-03-09_hardware-proxy-selection.md. State plainly
% whether the unoptimized model met a stated real-time threshold, since no
% quantization or other embedded optimization was applied -- this is a
% feasibility check, not an optimization result.

%%%%%%%%%%%%%%%%%%%%%%%%%%%%%%%%%%%%%%%%%%%%%%%%%
\section{Cross-Cutting Observations}\label{sec:results_general_observations}
% TODO: collects findings that don't belong to any single section above:
% training-time comparison across bands/setups; whether the synthetic-data
% practicability verdict (\Cref{sec:results_practicability_verdict}) is
% consistent with the band-comparison results
% (\Cref{sec:results_band_granularity}); an explicit closing statement that the
% knowledge-distillation and embedded-benchmark results
% (\Cref{sec:results_kd_and_embedded}) are preliminary/indicative given the time
% budget, setting up the Further Work section of Chapter
% \ref{chapter:discussion_conclusion} rather than being over-interpreted here.
